% !PRINTED 764
\weberpage{764}
4. \gesperrt{Der Beweis des Theorems I. ist hierdurch auf den Beweis des folgenden speciellen Falles zurückgeführt: Besteht eine Gleichung}

\weq{11}{C_1e^{z_1}+C_2e^{z_2}+\cdots+C_me^{z_m} = 0,}

\gesperrt{in der die $z_1, z_2, \ldots, z_m$ von einander verschiedene algebraische Zahlen sind, in der $C_1, C_2, \ldots, C_m$ und die Verbindungen}

\weq{12}{C_1g(z_1)+C_2g(z_2)+\cdots+C_mg(z_m)}

\gesperrt{rationale Zahlen sind, wenn $g(z)$ irgend eine ganze Function mit rationalen Coëfficienten ist, so müssen die $C_1, C_2, \ldots, C_m$ alle verschwinden.}

Um diesen Satz zu beweisen, bestimmen wir zunächst eine positive ganze rationale Zahl $c$ so, dass

\weq{13}{x_1 = cz_1,\ x_2 = cz_2,\ \ldots,\ x_m = cz_m}

ganze algebraische Zahlen werden (§.\,133, 5.). Wir nehmen die Coëfficienten $C_1, C_2, \ldots, C_m$ als \gesperrt{ganze} rationale Zahlen an, und bilden eine ganze Function

\weq{14}{\varphi(x) = ax^r+a_1x^{r-1}+\cdots+a_r}

mit ganzen rationalen Zahlencoëfficienten, die folgende Eigenschaften hat:

1) Die Zahlen $x_1, x_2, \ldots, x_m$ kommen unter den Wurzeln von $\varphi(x)=0$ vor.

2) Die Summe

\[
C_1\varphi'(x_1)+C_2\varphi'(x_2)+\cdots+C_m\varphi'(x_m) = k,
\]

die nach den Voraussetzungen (12), (13) eine ganze rationale Zahl ist, ist von Null verschieden.

Um einzusehen, dass eine solche Function $\varphi(x)$ immer existirt, nehme man zunächst eine Function $\chi(x)$, die nur der Bedingung 1) genügt, und der zweiten, dass keine der Zahlen $x_1, x_2, \ldots, x_m$ eine Doppelwurzel von $\chi(x)$ ist, eine solche Function existirt offenbar [man kann z. B., um eine Function $\chi$ zu bilden, die Norm des Productes $(x-x_1)(x-x_2)\cdots(x-x_m)$ von gemeinsamen Theilern mit ihren Derivirten befreien]. Setzt man dann

\[
\varphi(x) = x^h\chi(x),\ \varphi'(x_1) = x_1^h\chi'(x_1),\ \varphi'(x_2) = x_2^h\chi'(x_2)\ldots,
\]
%CONT>>
